\subsection{Randwertprobleme der Potentialtheorie}
Homogenes Randwertproblem: Beide Grenzen haben Potential 0.\\
Zu lösen ist die \textsc{Poisson}-Gleichung $\div(\varepsilon \nabla \Phi) = -\rho$ auf $\interior \Omega$:\\
\begin{tabular*}{\columnwidth}{@{\extracolsep\fill}llll@{}}
    Nr. & RWP & Randbedingungen auf $\partial \Omega$ & Lösung\\
    1. & Dirichlet & $\Phi\big|_{\partial\Omega} = \Phi_D$ & eindeutig $\Phi \in \mathcal C^2$\\[0.5em]
    2. & Neumann & $\frac{\partial \Phi}{\partial \vec n}\Big|_{\partial\Omega} = F_N$ & eindeutig $(\Phi + C) \in \mathcal C^2$\\[0.5em]
    3. & Gemischt & $\left( \Phi + k \frac{\partial \Phi}{\partial \vec n}\right)\Big|_{\partial\Omega} = F_N$ & eindeutig $\Phi \in \mathcal C^2$\\
\end{tabular*}
Mit Richtungsableitung $\frac{\partial \Phi}{\partial \vec n}\big|_{1/2} = \lim\limits_{\vec r - \vec r_0 \ra 0} \vec n(\vec r_0)  \nabla\Phi(\vec r)$ \quad $\vec r \in \Omega_{1/2}$\\
\\
Lösungsansatz: $\Phi = \Phi^{(0)} + \varphi$\\
$\Phi^{(0)}:$ erfüllt hom. DGL und inhom. RB\\
$\varphi:$ erfüllt inhom. DGL und hom. RB\\
\\
In den meisten Elektrostatischen Problemen gilt $\rho = 0$, da sich die Ladung nur auf den Grenzflächen von Leitern befindet und nicht im Gebiet $\Omega$ in dem die Lösung von $\Phi$ gesucht wird.\\
In der Praxis sind die meisten RWPs gemischt, wie Leiterkontakte oder Wärmeleitung\\

Mehrelektroden-Kondensator Q-RWP:\\
$\div(\varepsilon \nabla \Phi) = 0$ in $\interior \Omega$ und $\int_{\partial \Omega_l} \varepsilon \frac{\partial \Phi}{\partial \vec n} \diff \vec a = Q_l$ und  besitz bis auf eine additive Konstante eine eindeutige Lösung

\subsubsection{Spektralzerlegung}
Lösungsverfahren:\\
\begin{enumerate}
    \item Ansatz: $\Phi = \Phi^{(0)} + \varphi$\\
          Finde hinreichend glatte Funktion $\Phi^{(0)}$ welche inhomogene Randgleichungen erfüllt
    \item Finde Eigenfunktionen von $\varphi$: $f=-\div(\varepsilon \nabla \vec b_\nu) = \lambda_\nu \vec b_\nu$\\
          Es gilt $\lambda_\nu > 0$.
    \item Ansatz $\varphi(\vec r) = \sum_{\nu = 1}^\infty \alpha_\nu b_\nu(\vec r)$\\
          Bestimmung der Entwicklungskoeffizienten: $a_\nu = \frac{<b_\nu|f>}{\lambda_\nu} = \frac{1}{\lambda_\nu}\int_\Omega b_\nu * f dv$\\
    \item Spektraldarstellung: $G(\vec r, \vec r') = \sum_{\nu = 1}^\infty b_\nu (\vec r) \frac{1}{\lambda_\nu} b_\nu (\vec r')^*$\\
\end{enumerate}



\subsection{Greenfunktion $G$}
Def: Lösung des RWP mit hom. Randbed. und Störung $\rho(\vec r) = \delta(\vec r - \vec r')$ (Einheitspunktladung bei $\vec r'$)\\
Poissongleichung $\Delta\varphi = -\frac{\rho}{\varepsilon_0}$ wird durch das Coulomb-Integral gelöst.
Allg. Lösung: $\Phi(\vec r) = \varphi(\vec r)+\psi(\vec r) = \int_\Omega G(\vec r, \vec r')\rho(\vec r') \diff^3 \vec r' + \psi(\vec r)$\\
für $\varepsilon(\vec r) = \varepsilon$: $\psi(\vec r) = -\varepsilon\iint_{\partial V^{(D)}}\left[\dfrac{\partial G(\vec r, \vec r')}{\partial\vec n'}\Phi_D(\vec r')\right]\diff a'+\varepsilon\iint_{\partial V^{(N)}}\left[G(\vec r, \vec r')\dfrac{\partial\Phi_N(\vec r')}{\partial n'}\right]\diff a'$
\\
Beispiel Punktladung: $G_{\text{Vac}}(\vec r, \vec r') = \frac{1}{4 \pi \varepsilon} \frac{1}{\norm{\vec r - \vec r'}}$

\subsubsection{Spektralzerlegung mit Greenfunktion}
Problem: \boxed{- \lpo \varphi = \tilde{f}}

\begin{itemize}
    \item Sperationsansatz für die Eigenfunktionen: \\
          $b(\vec r) = b_1 (x_1) b_2(x_2) b_3(x_3)$
    \item $- \frac{b_1''(x_1)}{b_1 (x_1)} - \frac{b_2''(x_2)}{b_2 (x_2)}  - \frac{b_3''(x_3)}{b_3 (x_3)} = \lambda$
    \item Aufteilen des Problems: \\
          $- \frac{b_1''(x_1)}{b_1 (x_1)} = \lambda_1$ \\
          $- \frac{b_2''(x_2)}{b_2 (x_2)} = \lambda_2$ \\
          $- \frac{b_3''(x_3)}{b_3 (x_3)} = \lambda_3$
    \item Lösungsansatz für $b_1, b_2, b_3$: \\
          $b_j (x_j) = A_j \sin(k_j x_j) + B_j \cos (k_j x_j)$ mit $k_j = \sqrt{\lambda_j}$
    \item $\Ra B_j = 0$ und $k_j L_j = n_j \pi$
    \item Eigenfunktionen lauten: \\
          $b_j ( x_j ) = A_j \sin (n_j \frac{\pi}{L_j} x_j)$
    \item Normiere die Eigenfunktionen: \\
          $1 \overset{!}{=} \int \limits_0^{L_k} b_j (x_j)^2 \diff x_j$
\end{itemize}

Die Greenfunktion lautet nun: \\
$G(\vec r, \vec r') = \sum \limits_{n_1, n_2 , n_3 \in \mathbb N} b_{n_1 n_2 n_3} (\vec r ) \frac{1}{\lambda_{n_1}\lambda_{n_2}\lambda_{n_3}} b_{n_1 n_2 n_3} (\vec r' )$
%TODO S.52 Skript Lösungen der Laplace-Glg.
\subsection*{Spiegelladungsmethode}
Konstruktion eines Ersatzproblems durch Spiegelung der negierten Ladung an einer ebenen leitenden Randfläche\\
%TODO S.63 Bild
\begin{align*}
    G_{\text{Halb}}(\vec r, \vec r_0) = \frac{1}{4 \pi \varepsilon} \left( \frac{1}{\norm{\vec r - \vec r_0}} - \frac{1}{\norm{\vec r - \vec r_0^*}}\right)
\end{align*}
analog für Winkelräume. Eventuell müssen die gespiegelten Ladungen wieder gespiegelt werden (möglicherweise unendlich oft), bis sich alles ausgleicht.

\subsection*{Multipolentwicklung}
Coulomb-Integral: $\Phi(\vec r) = \int_{\R^3} G_{\ir vac}(\vec r, \vec r')\rho(\vec r') \diff^3 \vec r' = \frac{1}{4\pi\varepsilon_0}\int_{\R^3}  \frac{\rho(\vec r)}{\abs{\vec r - \vec r'}}\diff v'$\\
Vereinfachung der Integraldarstellung durch Taylorentwicklung des Integralkerns $\frac{1}{\abs{\vec r - \vec r'}}$ unter der Annahme $\abs{\vec r'} < \abs{\vec r}$:
$\Phi(\vec r) = \frac{1}{4\pi\varepsilon_0}\frac{1}{r}Q + \frac{1}{4\pi\varepsilon_0}\frac{\vec r  \vec p}{r^3} \mp \hdots$


\subsection{Stationäre Ströme und RWP}
Einflüsse: Drift, Diffusion, Hall-Effekt, Seebeck-Effekt\\
Drift-Diffusionsmodell:\\
\boxed{ \begin{array}{rll} \vec j = & \underset{\text{Driftstrom}}{\sum\limits_{\alpha = 1}^N |q_\alpha| n_\alpha \mu_\alpha \vec E} & - \underset{\text{Diffusionsstrom}}{\sum\limits_{\alpha = 1}^N q_\alpha D_\alpha \nabla n_\alpha} + \\[2em] + & \underset{\text{Halleffekt}}{\sum\limits_{\alpha = 1}^N \sigma_\alpha R_\alpha^H \vec \j_\alpha \times \vec B} & \underset{\text{Seebeck}}{-\sum\limits_{\alpha = 1}^N \sigma_\alpha P_\alpha \nabla T} \end{array}
}\\
$-\sum\limits_{\alpha = 1}^N \sigma_\alpha P_\alpha \nabla T$

%$\sum\limits_{\alpha = 1}^N \sigma_\alpha R_\alpha^H \vec \j_\alpha \times \vec B
